\label{apdx:call-for-collaborators}

\begin{center}
	\textit{The following recruitment text was prepared as a draft for identifying potential research collaborators. This document represents the intended approach to participant recruitment, subject to IRB approval.}
\end{center}

\vspace{1em}

\section*{Call for Collaborators: Exploring Sound, Culture, and Technology Together}

I'm working on a creative research project that brings together \textbf{sound, materials, movement, and technology} and I am looking for a small group of collaborators (about 5 people) to explore these ideas with me.

\subsection*{What is this about?}

I'm building sonic and computational tools that reflect specific cultural ways of knowing and making. Think: \textbf{algorithms inspired by cultural rhythms, interfaces that respond to embodied practice, sculptural forms that generate sound, and systems that carry memory and meaning.}

This project asks: \textit{What happens when we design technology from cultural epistemologies instead of just adapting to what already exists?}

\subsection*{What would collaboration look like?}

You'd be invited to:

\begin{itemize}[noitemsep]
	\item \textbf{Listen, respond, and critique} prototypes and experiments I'm working on
	\item \textbf{Share your perspectives} on sound, culture, technology, materials, and embodied practice
	\item \textbf{Engage in creative dialogue} about what feels authentic, what works, what doesn't
	\item \textbf{Co-create knowledge} through conversation and feedback sessions
\end{itemize}

Think of it as \textbf{creative critique sessions}: informal, generative, and conversational. \emph{Not interviews or data collection, but genuine collaboration.}

\subsection*{Who would be a good fit?}

You might be interested if you:

\begin{itemize}[noitemsep]
	\item Have experience with \textbf{music-making, sound design, sonic practices, or movement-based work} (any tradition)
	\item Work with \textbf{materials, physical objects, or embodied practices} (sculpture, crafts, dance, performance)
	\item Are curious about \textbf{cultural approaches to technology} or frustrated by tools that don't fit how you work
	\item Have thoughts about \textbf{rhythm, improvisation, gesture, call-and-response, or other cultural logics}
	\item Make things with your hands, body, code, sensors, or sound
	\item Care about \textbf{who gets to shape technology} and decolonial approaches to creativity
\end{itemize}

\textbf{You don't need to be an expert}. Just someone curious, reflective, and willing to think out loud.

\subsection*{What's in it for you?}

\begin{itemize}[noitemsep]
	\item Learn about creative coding, physical computing, and experimental sound systems
	\item Access to tools, techniques, and resources from the project
	\item Be part of shaping how cultural knowledge informs tech design
	\item Co-ownership of any collaborative knowledge or creations
	\item Community and connection with others exploring similar ideas
\end{itemize}

\subsection*{Time commitment?}

Flexible! Probably \textbf{3-5 sessions over several months} (1-2 hours each). We'll work around your schedule. Sessions can be in-person or virtual.

\subsection*{Interested?}

If this sounds exciting, reach out! I'd love to chat more about what this could look like and see if it's a good fit.

\textbf{Contact}: dkw34@drexel.edu

\vspace{1em}

\begin{center}
	\rule{0.8\textwidth}{0.5pt}
	
	\small\textit{This is part of my dissertation research at Drexel University, exploring how cultural epistemologies and speculative practices can guide the creation of new sonic and computational technologies.}
\end{center}